\section{Partitions induced by functions}

\import{set/partition.tex}
\import{function.tex}
\import{set/finite-sequence.tex}
\import{set/indexed-family.tex}

The fibres of a function are its canonical partition of the domain.  This is
the set-coded counterpart of quotienting by the kernel equivalence relation,
but it avoids introducing that relation when only the partition is needed.

\begin{definition}\label{function_fibre}
    $\fibre{f}{y}=\{x\in\dom{f}\mid f(x)=y\}$.
\end{definition}

\begin{proposition}\label{function_fibre_iff}
    Suppose $f$ is a function.
    Then $x\in\fibre{f}{y}$ iff $x\in\dom{f}$ and $f(x)=y$.
\end{proposition}
\begin{proof}
    Follows by \obv{function_fibre}.
\end{proof}

\begin{proposition}\label{function_fibre_restriction}
    Suppose $f$ is a function and $A\subseteq\dom{f}$.
    Then $\fibre{\restrl{f}{A}}{i}=\fibre{f}{i}\inter A$.
\end{proposition}
\begin{proof}
    Show that for every $x\in\fibre{\restrl{f}{A}}{i}$
        we have $x\in\fibre{f}{i}\inter A$.
    \begin{subproof}
        Fix $x\in\fibre{\restrl{f}{A}}{i}$.
        Hence $x\in\dom{\restrl{f}{A}}$ and
            $(\restrl{f}{A})(x)=i$ by \obv{function_fibre_iff}.
        Hence $x\in A$ by
            \auto{dom_of_restrl_eq_inter,inter_elim_right}.
        Hence $(\restrl{f}{A})(x)=f(x)$
            by \obv{restrl_of_function_apply}.
        Hence $x\in\fibre{f}{i}$ by \obv{function_fibre_iff,subseteq}.
        Follows by \obv{inter_intro}.
    \end{subproof}
    Show that for every $x\in\fibre{f}{i}\inter A$
        we have $x\in\fibre{\restrl{f}{A}}{i}$.
    \begin{subproof}
        Fix $x\in\fibre{f}{i}\inter A$.
        Hence $x\in\fibre{f}{i}$ and $x\in A$ by \obv{inter}.
        Hence $x\in\dom{f}$ and $f(x)=i$
            by \obv{function_fibre_iff}.
        Hence $x\in\dom{\restrl{f}{A}}$
            by \auto{dom_of_restrl_eq_inter,inter_intro}.
        $(\restrl{f}{A})(x)=f(x)$
            by \obv{restrl_of_function_apply}.
        Follows by \obv{function_fibre_iff}.
    \end{subproof}
    Follows by set extensionality.
\end{proof}

\begin{proposition}\label{function_fibre_nested_restriction}
    Suppose $f$ is a function and $A\subseteq\dom{f}$ and $B\subseteq A$.
    Then $\fibre{\restrl{f}{B}}{i}
        =\fibre{\restrl{f}{A}}{i}\inter B$.
\end{proposition}
\begin{proof}
    $B\subseteq\dom{f}$ by \obv{subseteq_transitive}.
    $\fibre{\restrl{f}{B}}{i}=\fibre{f}{i}\inter B$
        by \obv{function_fibre_restriction}.
    $\fibre{\restrl{f}{A}}{i}=\fibre{f}{i}\inter A$
        by \obv{function_fibre_restriction}.
    Follows by \auto{setext,inter,subseteq}.
\end{proof}

\begin{definition}\label{indexed_fibres}
    $\indexedFibres{f}{I}=\{(i,\fibre{f}{i})\mid i\in I\}$.
\end{definition}

\begin{proposition}\label{indexed_fibres_apply}
    Suppose $f$ is a function.
    Let $i\in I$.
    Then $(\indexedFibres{f}{I})(i)=\fibre{f}{i}$.
\end{proposition}
\begin{proof}
    $\indexedFibres{f}{I}$ is a relation
        by \auto{indexed_fibres,relation}.
    $\indexedFibres{f}{I}$ is right-unique
        by \auto{indexed_fibres,rightunique,pair_eq_iff}.
    $(i,\fibre{f}{i})\in\indexedFibres{f}{I}$
        by \obv{indexed_fibres}.
    Follows by \obv{function_apply_intro}.
\end{proof}

\begin{proposition}\label{indexed_fibres_type}
    Suppose $f$ is a function.
    Then $\indexedFibres{f}{I}\in\funs{I}{\pow{\dom{f}}}$.
\end{proposition}
\begin{proof}
    Let $P=\indexedFibres{f}{I}$.
    $P$ is a relation by \auto{indexed_fibres,relation}.
    $P$ is right-unique by \auto{indexed_fibres,rightunique,pair_eq_iff}.
    $\dom{P}=I$ by \auto{setext,dom_iff,indexed_fibres,pair_eq_iff}.
    Show that for every $i\in I$ we have $P(i)\in\pow{\dom{f}}$.
    \begin{subproof}
        Fix $i\in I$.
        $P(i)=\fibre{f}{i}$ by \auto{indexed_fibres_apply}.
        $\fibre{f}{i}\subseteq\dom{f}$
            by \obv{subseteq,function_fibre_iff}.
        Follows by \obv{powerset_intro}.
    \end{subproof}
    Hence $P$ is a function to $\pow{\dom{f}}$.
    Follows by \auto{funs_intro}.
\end{proof}

\begin{proposition}\label{indexed_fibres_pairwise_disjoint}
    Suppose $f$ is a function.
    Then $\indexedFibres{f}{I}$ is pairwise disjoint on $I$.
\end{proposition}
\begin{proof}
    Let $P=\indexedFibres{f}{I}$.
    Show that for all $i,j\in I$ such that $i\neq j$
        we have $P(i)$ is disjoint from $P(j)$.
    \begin{subproof}
        Fix $i,j\in I$.
        Assume $i\neq j$.
        $P(i)=\fibre{f}{i}$ and $P(j)=\fibre{f}{j}$
            by \auto{indexed_fibres_apply}.
        Suppose not.
        Take $x$ such that $x\in P(i),P(j)$ by \auto{disjoint}.
        Hence $f(x)=i$ and $f(x)=j$ by \auto{function_fibre_iff}.
        Hence $i=j$.
        Contradiction.
    \end{subproof}
    Follows by \auto{pairwise_disjoint_on}.
\end{proof}

\begin{proposition}\label{indexed_fibres_cover_values_in}
    Suppose $f$ is a function.
    Then $\unions{\ran{\indexedFibres{f}{I}}}
        =\{x\in\dom{f}\mid f(x)\in I\}$.
\end{proposition}
\begin{proof}
    Let $P=\indexedFibres{f}{I}$.
    $P\in\funs{I}{\pow{\dom{f}}}$ by \auto{indexed_fibres_type}.
    Hence $P$ is a function and $\dom{P}=I$
        by \obv{funs_is_function,funs}.
    Show that $\unions{\ran{P}}
        \subseteq\{y\in\dom{f}\mid f(y)\in I\}$.
    \begin{subproof}
        It suffices to show that for every $x\in\unions{\ran{P}}$
            we have $x\in\{y\in\dom{f}\mid f(y)\in I\}$
            by \obv{subseteq}.
        Fix $x\in\unions{\ran{P}}$.
        Take $B\in\ran{P}$ such that $x\in B$ by \obv{unions_iff}.
        Take $i\in\dom{P}$ such that $P(i)=B$ by \obv{fun_ran_iff}.
        Hence $i\in I$ by \obv{funs}.
        $P(i)=\fibre{f}{i}$ by \auto{indexed_fibres_apply}.
        Hence $x\in\dom{f}$ and $f(x)=i$
            by \auto{function_fibre_iff}.
        Follows by \auto.
    \end{subproof}
    Show that $\{y\in\dom{f}\mid f(y)\in I\}
        \subseteq\unions{\ran{P}}$.
    \begin{subproof}
        It suffices to show that for every
            $x\in\{y\in\dom{f}\mid f(y)\in I\}$
            we have $x\in\unions{\ran{P}}$ by \obv{subseteq}.
        Fix $x\in\{y\in\dom{f}\mid f(y)\in I\}$.
        Hence $x\in\dom{f}$ and $f(x)\in I$ by \auto.
        Hence $x\in\fibre{f}{f(x)}$ by \obv{function_fibre_iff}.
        $P(f(x))=\fibre{f}{f(x)}$ by \auto{indexed_fibres_apply}.
        $P(f(x))\in\ran{P}$ by \auto{funs,fun_ran_iff}.
        Follows by \auto{unions_iff}.
    \end{subproof}
    Follows by \obv{subseteq_antisymmetric}.
\end{proof}

\begin{proposition}\label{indexed_fibres_cover_domain}
    Let $f\in\funs{D}{I}$.
    Then $\unions{\ran{\indexedFibres{f}{I}}}=D$.
\end{proposition}
\begin{proof}
    $f$ is a function and $\dom{f}=D$
        by \obv{funs_is_function,funs}.
    $\unions{\ran{\indexedFibres{f}{I}}}
        =\{x\in\dom{f}\mid f(x)\in I\}$
        by \obv{indexed_fibres_cover_values_in}.
    Show that $\{x\in\dom{f}\mid f(x)\in I\}=\dom{f}$.
    \begin{subproof}
        Show that $\{x\in\dom{f}\mid f(x)\in I\}\subseteq\dom{f}$.
        \begin{subproof}
            Follows by \auto{subseteq}.
        \end{subproof}
        Show that $\dom{f}\subseteq\{x\in\dom{f}\mid f(x)\in I\}$.
        \begin{subproof}
            It suffices to show that for every $x\in\dom{f}$
                we have $x\in\{y\in\dom{f}\mid f(y)\in I\}$
                by \obv{subseteq}.
            Fix $x\in\dom{f}$.
            Hence $x\in D$.
            $f(x)\in I$ by \obv{funs_type_apply}.
            Follows by \auto.
        \end{subproof}
        Follows by \obv{subseteq_antisymmetric}.
    \end{subproof}
\end{proof}

\begin{definition}\label{function_fibres}
    $\fibres{f}=\{\fibre{f}{y}\mid y\in\ran{f}\}$.
\end{definition}

\begin{proposition}\label{function_fibres_elim}
    Suppose $B\in\fibres{f}$.
    Then there exists $y\in\ran{f}$ such that $B=\fibre{f}{y}$.
\end{proposition}
\begin{proof}
    Follows by \obv{function_fibres}.
\end{proof}

\begin{proposition}\label{function_fibres_intro}
    Suppose $y\in\ran{f}$.
    Then $\fibre{f}{y}\in\fibres{f}$.
\end{proposition}
\begin{proof}
    Follows by \obv{function_fibres}.
\end{proof}

\begin{proposition}\label{function_fibres_nonempty}
    Suppose $f$ is a function.
    Then $\emptyset\notin\fibres{f}$.
\end{proposition}
\begin{proof}
    Suppose not.
    Take $y\in\ran{f}$ such that $\emptyset=\fibre{f}{y}$
        by \obv{function_fibres_elim}.
    Take $x\in\dom{f}$ such that $f(x)=y$ by \obv{fun_ran_iff}.
    Hence $x\in\fibre{f}{y}$ by \obv{function_fibre_iff}.
    Contradiction by \obv{emptyset}.
\end{proof}

\begin{proposition}\label{distinct_function_fibres_disjoint}
    Suppose $f$ is a function.
    Suppose $B,C\in\fibres{f}$ and $B\neq C$.
    Then $B$ is disjoint from $C$.
\end{proposition}
\begin{proof}
    Take $y\in\ran{f}$ such that $B=\fibre{f}{y}$
        by \obv{function_fibres_elim}.
    Take $z\in\ran{f}$ such that $C=\fibre{f}{z}$
        by \obv{function_fibres_elim}.
    Suppose not.
    Take $x$ such that $x\in B,C$ by \auto{disjoint}.
    Then $f(x)=y$ and $f(x)=z$ by \auto{function_fibre_iff}.
    Hence $y=z$.
    Thus $B=C$.
    Contradiction.
\end{proof}

\begin{proposition}\label{function_fibres_cover_domain}
    Suppose $f$ is a function.
    Then $\unions{\fibres{f}}=\dom{f}$.
\end{proposition}
\begin{proof}
    Show that $\unions{\fibres{f}}\subseteq\dom{f}$.
    \begin{subproof}
        It suffices to show that every element of $\unions{\fibres{f}}$
            is an element of $\dom{f}$ by \obv{subseteq}.
        Fix $x\in\unions{\fibres{f}}$.
        Take $B\in\fibres{f}$ such that $x\in B$ by \obv{unions_iff}.
        Take $y\in\ran{f}$ such that $B=\fibre{f}{y}$
            by \obv{function_fibres_elim}.
        Follows by \auto{function_fibre_iff}.
    \end{subproof}
    Show that $\dom{f}\subseteq\unions{\fibres{f}}$.
    \begin{subproof}
        It suffices to show that every element of $\dom{f}$
            is an element of $\unions{\fibres{f}}$ by \obv{subseteq}.
        Fix $x\in\dom{f}$.
        $(x,f(x))\in f$ by \obv{function_apply_elim}.
        Hence $f(x)\in\ran{f}$ by \obv{ran_intro}.
        Thus $\fibre{f}{f(x)}\in\fibres{f}$
            by \obv{function_fibres_intro}.
        $x\in\fibre{f}{f(x)}$ by \obv{function_fibre_iff}.
        Follows by \auto{unions_iff}.
    \end{subproof}
    Follows by \obv{subseteq_antisymmetric}.
\end{proof}

\begin{theorem}\label{function_fibres_partition}
    Suppose $f$ is a function.
    Then $\fibres{f}$ is a partition of $\dom{f}$.
\end{theorem}
\begin{proof}
    $\emptyset\notin\fibres{f}$ by \obv{function_fibres_nonempty}.
    For all $B,C\in\fibres{f}$ such that $B\neq C$
        we have $B$ is disjoint from $C$
        by \obv{distinct_function_fibres_disjoint}.
    Thus $\fibres{f}$ is a partition by \obv{partition}.
    Follows by \obv{function_fibres_cover_domain}.
\end{proof}

\begin{corollary}\label{sequence_fibres_partition}
    Suppose $s$ is a finite sequence in $A$.
    Then $\fibres{s}$ is a partition of $\dom{s}$.
\end{corollary}
\begin{proof}
    $s$ is a function by \auto{finite_sequence_domain,funs_is_function}.
    Follows by \obv{function_fibres_partition}.
\end{proof}
